\documentclass[11pt]{article}
\usepackage[margin=0.8in]{geometry}
\usepackage{amsmath,amssymb}
\usepackage{wasysym}
\usepackage{array}
\setlength{\parskip}{6pt}
\begin{document}

\begin{center}
{\Large CSC343 Assignment 3 -- Part 2 Solutions}
\end{center}

%========================================================================
\section*{1. Relation $R_1(L,M,N,O,P,Q,R,S)$}
Given FDs:
\[
S_1 = \{L \to NQ,\; MNR \to O,\; O \to M,\; NQ \to LS,\; S \to OPR\}
\]

%------------------------------------------------------------------------
\subsection*{(a) BCNF violations}

An FD $X \to Y$ violates BCNF iff $X$ is not a superkey of $R_1$.
We compute the closure of each LHS in $S_1$:
\begin{itemize}
\item $L^+ = \{L\}
  \xrightarrow{L \to NQ} \{L,N,Q\}
  \xrightarrow{NQ \to LS} \{L,N,Q,S\}
  \xrightarrow{S \to OPR} \{L,N,O,P,Q,R,S\}$\\
  $\xrightarrow{O \to M} \{L,M,N,O,P,Q,R,S\}$.
  $L$ is a superkey.
\item $(NQ)^+ = \{N,Q\}
  \xrightarrow{NQ \to LS} \{L,N,Q,S\}
  \xrightarrow{S \to OPR} \{L,N,O,P,Q,R,S\}$\\
  $\xrightarrow{O \to M} \{L,M,N,O,P,Q,R,S\}$.
  $NQ$ is a superkey.
\item $(MNR)^+ = \{M,N,R\}
  \xrightarrow{MNR \to O} \{M,N,O,R\}
  \xrightarrow{O \to M} \{M,N,O,R\}$.
  Not a superkey.
\item $O^+ = \{O\}
  \xrightarrow{O \to M} \{M,O\}$.
  Not a superkey.
\item $S^+ = \{S\}
  \xrightarrow{S \to OPR} \{O,P,R,S\}
  \xrightarrow{O \to M} \{M,O,P,R,S\}$.
  Not a superkey.
\end{itemize}

BCNF-violating FDs (LHS is not a superkey):
\[
MNR \to O,\quad O \to M,\quad S \to OPR.
\]

%------------------------------------------------------------------------
\subsection*{(b) BCNF decomposition}

\textbf{Step 1.} Start with $R_1(LMNOPQRS)$.
Choose violating FD $O \to M$ (where $O^+ = \{M,O\}$).
\begin{itemize}
\item $R_{11}(MO)$ \quad (attributes in $O^+$)
\item $R_{12}(LNOPQRS)$ \quad ($R_1$ attributes minus $M$, keeping $O$)
\end{itemize}

Projected FDs on $R_{12}(LNOPQRS)$:
$\{L \to NQ,\; NQ \to LS,\; S \to OPR\}$.

($MNR \to O$ and $O \to M$ cannot project onto $R_{12}$ since $M \notin R_{12}$.)

Check BCNF in $R_{12}$: $L$ and $NQ$ are still superkeys (same closures
as before, restricted to $R_{12}$). But
$S^+ = \{S\} \to \{O,P,R,S\}$, so $S$ is not a superkey.
Thus $S \to OPR$ violates BCNF in $R_{12}$.

\medskip
\textbf{Step 2.} Decompose $R_{12}$ using $S \to OPR$
(where $S^+$ restricted to $R_{12}$ is $\{O,P,R,S\}$):
\begin{itemize}
\item $R_{121}(OPRS)$
\item $R_{122}(LNQS)$
\end{itemize}

Verify each final relation is in BCNF:
\begin{itemize}
\item $R_A(LNQS)$: $L^+ = \{L,N,Q,S\} = R_A$ and
      $(NQ)^+ = \{L,N,Q,S\} = R_A$. Both LHS are superkeys. BCNF.
\item $R_B(MO)$: $O^+ = \{M,O\} = R_B$. Superkey. BCNF.
\item $R_C(OPRS)$: $S^+ = \{O,P,R,S\} = R_C$. Superkey. BCNF.
\end{itemize}

\textbf{Final BCNF decomposition}
(relations in alphabetical order):
\[
R_A(LNQS),\quad R_B(MO),\quad R_C(OPRS).
\]

Projected FDs:
\begin{itemize}
\item On $R_A(LNQS)$: $L \to NQ,\; NQ \to LS$.
\item On $R_B(MO)$: $O \to M$.
\item On $R_C(OPRS)$: $S \to OPR$.
\end{itemize}

%------------------------------------------------------------------------
\subsection*{(c) Dependency preservation}

The union of all projected FDs is:
$\{L \to NQ,\; NQ \to LS,\; O \to M,\; S \to OPR\}$.

We check whether $MNR \to O$ can be derived from these.
Compute $(MNR)^+$ under the projected FDs:
\[
\{M,N,R\}
  \;\xrightarrow{\text{no FD fires}}\;
  \{M,N,R\}.
\]
$M$, $N$, and $R$ never appear together in any single decomposed relation,
so no projected FD has its LHS contained in $\{M,N,R\}$.
Since $O \notin \{M,N,R\}$, the FD $MNR \to O$ is \textbf{lost}.

The decomposition is \textbf{not dependency-preserving}.

%------------------------------------------------------------------------
\subsection*{(d) Chase test for lossless join}

Decomposition: $\{R_B(MO),\; R_C(OPRS),\; R_A(LNQS)\}$.

Assign a lowercase letter to each attribute
($l, m, n, o, p, q, r, s$ for $L, M, N, O, P, Q, R, S$).
Each row gets letters for the attributes in its relation and
arbitrary numbers elsewhere.

\medskip
\textbf{Initial tableau:}

\begin{center}
\begin{tabular}{|l|c c c c c c c c|}
\hline
       & $L$ & $M$ & $N$ & $O$ & $P$ & $Q$ & $R$ & $S$ \\
\hline
$MO$   & 1   & $m$ & 2   & $o$ & 3   & 4   & 5   & 6   \\
$OPRS$ & 7   & 8   & 9   & $o$ & $p$ & 10  & $r$ & $s$ \\
$LNQS$ & $l$ & 11  & $n$ & 12  & 13  & $q$ & 14  & $s$ \\
\hline
\end{tabular}
\end{center}

\textbf{Apply $S \to OPR$:}
Rows $OPRS$ and $LNQS$ agree on $S = s$, so $O, P, R$ must match.
Copy from $OPRS$: in row $LNQS$, replace $12 \to o$,
$13 \to p$, $14 \to r$.

\begin{center}
\begin{tabular}{|l|c c c c c c c c|}
\hline
       & $L$ & $M$ & $N$ & $O$          & $P$          & $Q$
       & $R$          & $S$ \\
\hline
$MO$   & 1   & $m$ & 2   & $o$          & 3            & 4
       & 5            & 6   \\
$OPRS$ & 7   & 8   & 9   & $o$          & $p$          & 10
       & $r$          & $s$ \\
$LNQS$ & $l$ & 11  & $n$ & $\mathbf{o}$ & $\mathbf{p}$ & $q$
       & $\mathbf{r}$ & $s$ \\
\hline
\end{tabular}
\end{center}

\textbf{Apply $O \to M$:}
All three rows agree on $O = o$, so $M$ must match.
Copy from $MO$: in row $OPRS$, replace $8 \to m$;
in row $LNQS$, replace $11 \to m$.

\begin{center}
\begin{tabular}{|l|c c c c c c c c|}
\hline
       & $L$ & $M$          & $N$ & $O$ & $P$ & $Q$ & $R$ & $S$ \\
\hline
$MO$   & 1   & $m$          & 2   & $o$ & 3   & 4   & 5   & 6   \\
$OPRS$ & 7   & $\mathbf{m}$ & 9   & $o$ & $p$ & 10  & $r$ & $s$ \\
$LNQS$ & $l$ & $\mathbf{m}$ & $n$ & $o$ & $p$ & $q$ & $r$ & $s$ \\
\hline
\end{tabular}
\end{center}

Row $LNQS$ now has a letter for every attribute:
$l, m, n, o, p, q, r, s$.
No further FD applications produce new changes.

Chase test \textbf{passed}. The decomposition is \textbf{lossless}.

\newpage
%========================================================================
\section*{2. Relation $R_2(A,B,C,D,E,F,G,H)$}
Given FDs:
\[
S_2 = \{AB \to C,\; C \to ABD,\; CFD \to E,\;
        E \to B,\; BF \to EC,\; B \to DA\}
\]

%------------------------------------------------------------------------
\subsection*{(a) Minimal basis}

\textbf{Step 1: Split RHS} into single-attribute FDs.
\[
\{\;AB \to C,\;
C \to A,\; C \to B,\; C \to D,\;
CDF \to E,\;
E \to B,\;
BF \to C,\; BF \to E,\;
B \to A,\; B \to D
\;\}
\]

\bigskip
\textbf{Step 2: Remove extraneous LHS attributes.}

Only FDs with $|\text{LHS}| > 1$: $AB \to C$, $CDF \to E$,
$BF \to C$, $BF \to E$.
\begin{itemize}
\item $AB \to C$: try removing $A$.
  $B^+ = \{B\}
   \xrightarrow{B \to D} \{B,D\}
   \xrightarrow{B \to A} \{A,B,D\}
   \xrightarrow{AB \to C} \{A,B,C,D\}$.
  $C \in B^+$, so $A$ is extraneous. Replace with $B \to C$.
\item $CDF \to E$: try removing $C$.
  $(DF)^+ = \{D,F\}$. $E \notin$. Keep $C$.\\
  Try removing $D$.
  $(CF)^+ = \{C,F\}
   \xrightarrow{C \to B} \{B,C,F\}
   \xrightarrow{B \to D} \{B,C,D,F\}
   \xrightarrow{BF \to E} \{B,C,D,E,F\}$.
  $E \in (CF)^+$, so $D$ is extraneous. Replace with $CF \to E$.\\
  Re-check $CF \to E$: $F^+ = \{F\}$, $E \notin$; keep $C$.
  $C^+ = \{C\} \to \{B,C\} \to \{A,B,C,D\}$, $E \notin$; keep $F$.
\item $BF \to C$: try removing $F$.
  $B^+ = \{B\} \xrightarrow{B \to C} \{B,C\}$.
  $C \in B^+$, so $F$ is extraneous. This gives duplicate $B \to C$;
  remove this FD.
\item $BF \to E$: $F^+ = \{F\}$, $E \notin$; keep $B$.
  $B^+ = \{B\} \to \{B,C\} \to \{A,B,C,D\}$, $E \notin$; keep $F$.
\end{itemize}

After Step 2:
$\{B \to C,\; C \to A,\; C \to B,\; C \to D,\; CF \to E,\;
  E \to B,\; BF \to E,\; B \to A,\; B \to D\}$

\bigskip
\textbf{Step 3: Remove redundant FDs.}

Remove each FD one at a time; compute closure of its LHS using the
\emph{remaining} FDs. Once an FD is discarded, it stays removed.

\begin{center}
\begin{tabular}{|c|l|p{6.8cm}|l|}
\hline
\# & FD & LHS closure \emph{without} this FD & Decision \\
\hline
1 & $B \to C$
& $B^+ = \{B\}
   \xrightarrow{B \to D} \{B,D\}
   \xrightarrow{B \to A} \{A,B,D\}$.
  $C \notin B^+$.
& keep \\
\hline
2 & $C \to A$
& $C^+ = \{C\}
   \xrightarrow{C \to B} \{B,C\}
   \xrightarrow{B \to A} \{A,B,C\}$.
  $A \in C^+$.
& \textbf{discard} \\
\hline
3 & $C \to D$
& (after discarding $C \to A$)
  $C^+ = \{C\}
   \xrightarrow{C \to B} \{B,C\}
   \xrightarrow{B \to D} \{B,C,D\}$.
  $D \in C^+$.
& \textbf{discard} \\
\hline
4 & $CF \to E$
& (after discarding $C \to A$, $C \to D$)
  $(CF)^+ = \{C,F\}
   \xrightarrow{C \to B} \{B,C,F\}
   \xrightarrow{BF \to E} \{B,C,E,F\}$.
  $E \in (CF)^+$.
& \textbf{discard} \\
\hline
5 & $C \to B$
& $C^+ = \{C\}$. No FD fires. $B \notin C^+$.
& keep \\
\hline
6 & $E \to B$
& $E^+ = \{E\}$. No FD fires. $B \notin E^+$.
& keep \\
\hline
7 & $BF \to E$
& $(BF)^+ = \{B,F\}
   \xrightarrow{B \to C} \{B,C,F\}
   \xrightarrow{B \to D} \{B,C,D,F\}
   \xrightarrow{B \to A} \{A,B,C,D,F\}$.
  $E \notin (BF)^+$.
& keep \\
\hline
8 & $B \to A$
& $B^+ = \{B\}
   \xrightarrow{B \to C} \{B,C\}
   \xrightarrow{B \to D} \{B,C,D\}$.
  $A \notin B^+$.
& keep \\
\hline
9 & $B \to D$
& $B^+ = \{B\}
   \xrightarrow{B \to C} \{B,C\}
   \xrightarrow{B \to A} \{A,B,C\}$.
  $D \notin B^+$.
& keep \\
\hline
\end{tabular}
\end{center}

\textbf{Minimal basis} (in alphabetical order):
\[
\boxed{\;B \to A,\quad B \to C,\quad B \to D,\quad
        BF \to E,\quad C \to B,\quad E \to B\;}
\]

%------------------------------------------------------------------------
\subsection*{(b) All keys for $R_2$}

Identify attributes by their LHS/RHS appearances in the minimal basis:

\begin{center}
\begin{tabular}{|c|c|c|l|}
\hline
Attribute & Appears on LHS & Appears on RHS & Conclusion \\
\hline
$A$ & --       & \checked & not necessarily in any key \\
$B$ & \checked & \checked & must check \\
$C$ & \checked & \checked & must check \\
$D$ & --       & \checked & not necessarily in any key \\
$E$ & \checked & \checked & must check \\
$F$ & \checked & --       & must be in every key \\
$G$ & --       & --       & must be in every key \\
$H$ & --       & --       & must be in every key \\
\hline
\end{tabular}
\end{center}

$F, G, H$ must be in every key.
$(FGH)^+ = \{F,G,H\} \neq R_2$, so we need at least one of $\{B,C,E\}$.
\begin{itemize}
\item $(BFGH)^+ = \{B,F,G,H\}
   \xrightarrow{B \to A} \{A,B,F,G,H\}
   \xrightarrow{B \to C} \{A,B,C,F,G,H\}
   \xrightarrow{B \to D} \{A,B,C,D,F,G,H\}
   \xrightarrow{BF \to E} \{A,B,C,D,E,F,G,H\}$. Key.
\item $(CFGH)^+ = \{C,F,G,H\}
   \xrightarrow{C \to B} \{B,C,F,G,H\}
   \xrightarrow{B \to A} \{A,B,C,F,G,H\}
   \xrightarrow{B \to D} \{A,B,C,D,F,G,H\}
   \xrightarrow{BF \to E} \{A,B,C,D,E,F,G,H\}$. Key.
\item $(EFGH)^+ = \{E,F,G,H\}
   \xrightarrow{E \to B} \{B,E,F,G,H\}
   \xrightarrow{B \to A} \{A,B,E,F,G,H\}
   \xrightarrow{B \to C} \{A,B,C,E,F,G,H\}
   \xrightarrow{B \to D} \{A,B,C,D,E,F,G,H\}$. Key.
\end{itemize}
$A, D$ are not on any LHS, so $AFGH$ and $DFGH$ cannot be keys.
Each of $BFGH$, $CFGH$, $EFGH$ is minimal since $FGH$ alone is not
a superkey. All candidate keys:
\[
\boxed{\;BFGH,\quad CFGH,\quad EFGH\;}
\]

%------------------------------------------------------------------------
\subsection*{(c) 3NF synthesis}

From the minimal basis, combine FDs with the same LHS:
\[
B \to ACD,\quad BF \to E,\quad C \to B,\quad E \to B.
\]

Create one relation per group (LHS $\cup$ RHS):
\begin{itemize}
\item From $B \to ACD$: $R_1(ABCD)$, key $B$.
\item From $BF \to E$: $R_2(BEF)$, key $BF$.
\item From $C \to B$: $R_3(BC)$, key $C$.
\item From $E \to B$: $R_4(BE)$, key $E$.
\end{itemize}

Remove subset relations:
\begin{itemize}
\item $R_3(BC) \subseteq R_1(ABCD)$: $\{B,C\} \subset \{A,B,C,D\}$.
      Remove $R_3$.
\item $R_4(BE) \subseteq R_2(BEF)$: $\{B,E\} \subset \{B,E,F\}$.
      Remove $R_4$.
\end{itemize}

Check if any remaining relation contains a candidate key for $R_2$:
\begin{itemize}
\item $R_1(ABCD)$: does not contain $F, G, H$. No.
\item $R_2(BEF)$: does not contain $G, H$. No.
\end{itemize}

Neither contains a key, so add a key relation: $R_5(BFGH)$.

\medskip
\textbf{Final 3NF decomposition} (alphabetical):
\[
\boxed{\;R_1(ABCD),\quad R_2(BEF),\quad R_3(BFGH)\;}
\]

\begin{itemize}
\item \textbf{Dependency-preserving}: by construction, every FD in the
      minimal basis is contained within some relation.
      $B \to ACD$ in $R_1$; $BF \to E$ in $R_2$;
      $C \to B$ in $R_1$ ($C, B \in ABCD$);
      $E \to B$ in $R_2$ ($E, B \in BEF$).
\item \textbf{Lossless}: $R_3(BFGH)$ contains candidate key $BFGH$.
\end{itemize}

%------------------------------------------------------------------------
\subsection*{(d) Redundancy}

Yes, redundancy is possible.

Consider $R_2(BEF)$. Its projected FDs are $BF \to E$ and $E \to B$.
The keys of $R_2$ are $BF$ and $EF$ (since $(EF)^+ = \{E,F\}
\xrightarrow{E \to B} \{B,E,F\} = R_2$).

The FD $E \to B$ violates BCNF in $R_2$ because $E$ is not a superkey
($(E)^+ = \{B,E\} \neq R_2$). Although $B$ is a prime attribute
(so 3NF is not violated), the BCNF violation means redundancy can occur.

Concrete example: if $e_1$ determines $b_1$, then tuples
$(b_1, e_1, f_1)$ and $(b_1, e_1, f_2)$ repeat the fact
``$E = e_1 \Rightarrow B = b_1$'' redundantly.
This also creates the risk of update anomalies.

\end{document}
